\scanpage{57}
\mspage{6}

% slanted
\margin{\uline{sans} axiome \textcircled{4} C'est la première fois qu'on a besoin de $W(2)$ sous sa forme forte, à l'exclusion de $W(2\,\mathrm{bis})$.}

\uline{Proposition 3} Soient $X,Y$ dans $\mathrm{Cat}_{/S}$, $f:X\unsure{\text{---}|\longrightarrow}Y$ une $S$-flèche. Supposons \inserted{$X,Y$ $W$-fibrés} que $\forall s\in S$, $f_s:X_s\longrightarrow Y_s$ soit \deleted{dans} dans $W$, alors $f\in W$.

Il s'impose d'utiliser l'axiome de localisation, donc regarder les $f_{/s}$ induits par $f$
\diagram{\begin{tikzcd}[ampersand replacement=\&,row sep=small,column sep=large]
X_{/s}\arrow[r,"f_{/s}"'] \& Y_{/s}\\
X_s\arrow[u,"i_s^X"]\arrow[r,"f_s"] \& Y_s\arrow[u,"i_s^Y"']
\end{tikzcd}}

\deleted{\unsure{dans les}} condition $a_1$ suffit à impliquer que \deleted{\unsure{les}} $i_s^X$ et $i_s^Y$ sont dans $W$, donc $f_{/s}$ est dans $W$ car $f_s$ l'est. (On n'a pas eu à utiliser l'axiome \textcircled{4}, et il suffisait de supposer que \deleted{\unsure{$f$}}$X,Y$ soient fibrés ``$W$-cohomologiques'' \deleted{\unsure{sur $S$}} (ou \inserted{$W$-fibrés} à g.) sur $S$, ou d'ailleurs, des $W$-fibrés à dr.

De plus :

\uline{Corollaire} \deleted{\unsure{$S$}} Soit $S_0\subset\mathrm{Ob}\,S$ une partie de $S$ telle que $S_0\longrightarrow\pi_0(S)$ surjectif. Alors dans la prop 3 il suffit de supposer $f_s\in W_{\unsure{4|s}}$ pour $s\in S_0$.

% slanted
\margin{sans axiome \textcircled{4}}

Il suffit de voir que si $u:s\longrightarrow s'$ est une flèche dans $\unsure{s|S}$, alors $f_s\in W\Longleftrightarrow f_{s'}\in W$. On utilise le fait que $f_s\in W\Longleftrightarrow f_{/s}\in W$, et de même pour $f_{s'}$. Donc à voir que $f_{/s}\in W\Longleftrightarrow f_{/s'}\in W$. Or on a

\diagram{\begin{tikzcd}[ampersand replacement=\&,row sep=large,column sep=large]
X_{/s}\arrow[r,"f_{/s}"]\arrow[d,"X_{/u}"'] \& Y_{/s}\arrow[d,"X_{/u}"]\\
X_{/s'}\arrow[r,"f_{/s'}"] \& Y_{/s'}
\end{tikzcd}}
où $X_{/u}$ et $Y_{/u}$ sont dans $W$, d'où le résultat.

Mais revenons à la réciproque : si $f\in W$, alors a-t-on $f_s\in W_s$ ? Il y a un argument dans ce sens dans $\overline{\underline{\mathrm{VI}}}$ p. 23, que je vais reprendre. L'idée \deleted{c'est de} est celle-ci. On considère la comp. conn. $Y_0$ de $Y$ contenant $s$, et on \deleted{\unsure{est}} construit un

% archive sheet number: 57
